\chapter{Данные}
\label{ch:data}

\section{Расширенное и детализированное описание данных и их использование в проекте онлайн-кинотеатра}

Проект онлайн-кинотеатра представляет собой комплексное решение, которое включает в себя бэкенд для хранения и управления данными, а также рекомендательную систему на основе машинного обучения. В качестве основного источника данных используется датасет \textbf{MovieLens 100K}, который является одним из самых популярных и широко применяемых наборов данных для задач рекомендательных систем. Этот датасет содержит информацию о фильмах, пользователях и их взаимодействиях (оценках). Ниже приведено детальное описание структуры данных, их обработки, интеграции в проект и дальнейшего использования.

\section{Структура данных в датасете MovieLens 100K}

Датасет MovieLens 100K состоит из нескольких файлов, каждый из которых содержит определённый тип данных. Эти данные были тщательно структурированы и интегрированы в проект для обеспечения максимальной эффективности и удобства использования.

\subsection*{Фильмы (Movies)}
Файл: \texttt{u.item}

Содержит информацию о фильмах, включая:
\begin{itemize}
    \item \textbf{ID фильма} (уникальный идентификатор, который используется для связи с другими таблицами).
    \item \textbf{Название фильма} (включая год выпуска, что позволяет легко идентифицировать фильм).
    \item \textbf{Дата выпуска} (в формате DD-MMM-YYYY, что позволяет анализировать временные тренды).
    \item \textbf{URL-ссылка на IMDb} (для получения дополнительной информации о фильме).
    \item \textbf{Жанры} (бинарные флаги для 19 жанров, таких как Action, Comedy, Drama, Horror, Romance и т.д.).
\end{itemize}

Пример строки:
\begin{verbatim}
1|Toy Story (1995)|01-Jan-1995|
|http://us.imdb.com/M/title-exact
?Toy%20Story%20(1995)|0|0|0|1|1|
1|0|0|0|0|0|0|0|0|0|0|0|0|0
\end{verbatim}

\subsection*{Пользователи (Users)}
Файл: \texttt{u.user}

Содержит информацию о пользователях:
\begin{itemize}
    \item \textbf{ID пользователя} (уникальный идентификатор, который используется для связи с таблицей оценок).
    \item \textbf{Возраст} (числовое значение, которое позволяет сегментировать пользователей по возрастным группам).
    \item \textbf{Пол} (категориальная переменная: M — мужской, F — женский).
    \item \textbf{Профессия} (категориальная переменная, например, educator, engineer, student, artist и т.д.).
    \item \textbf{Почтовый индекс} (географическая информация, которая может быть использована для анализа региональных предпочтений).
\end{itemize}

Пример строки:
\begin{verbatim}
1|24|M|technician|85711
\end{verbatim}

\subsection*{Оценки (Ratings)}
Файл: \texttt{u.data}

Содержит информацию о том, как пользователи оценивали фильмы:
\begin{itemize}
    \item \textbf{ID пользователя} (связь с таблицей пользователей).
    \item \textbf{ID фильма} (связь с таблицей фильмов).
    \item \textbf{Оценка} (по шкале от 1 до 5, где 1 — низкая оценка, 5 — высокая оценка).
    \item \textbf{Временная метка} (timestamp, который позволяет анализировать временные закономерности в поведении пользователей).
\end{itemize}

Пример строки:
\begin{verbatim}
196 242 3 881250949
\end{verbatim}

\subsection*{Дополнительные данные}
\begin{itemize}
    \item \textbf{Информация о жанрах}: В датасете используется бинарное кодирование жанров, что позволяет легко фильтровать фильмы по категориям.
    \item \textbf{Временные метки}: Позволяют анализировать временные закономерности в поведении пользователей.
\end{itemize}

\section{Интеграция данных в проект}

\subsection*{Хранение данных в бэкенде}
Данные из датасета MovieLens 100K были интегрированы в базу данных проекта. Для этого использовалась следующая структура:

\begin{itemize}
    \item \textbf{Таблица \texttt{movies}}:
        \begin{itemize}
            \item \texttt{id} (уникальный идентификатор фильма).
            \item \texttt{title} (название фильма).
            \item \texttt{release\_date} (дата выпуска).
            \item \texttt{genres} (список жанров в формате JSON или отдельная таблица для связи многие-ко-многим).
            \item \texttt{imdb\_url} (ссылка на IMDb).
        \end{itemize}
    \item \textbf{Таблица \texttt{users}}:
        \begin{itemize}
            \item \texttt{id} (уникальный идентификатор пользователя).
            \item \texttt{age} (возраст).
            \item \texttt{gender} (пол).
            \item \texttt{occupation} (профессия).
            \item \texttt{zip\_code} (почтовый индекс).
        \end{itemize}
    \item \textbf{Таблица \texttt{ratings}}:
        \begin{itemize}
            \item \texttt{user\_id} (ID пользователя).
            \item \texttt{movie\_id} (ID фильма).
            \item \texttt{rating} (оценка от 1 до 5).
            \item \texttt{timestamp} (временная метка).
        \end{itemize}
\end{itemize}

\section{Расширенное и детализированное}

Пример строки:
\begin{verbatim}
196 242 3 881250949
\end{verbatim}

\section{Дополнительные данные}
\subsection*{Информация о жанрах}
В датасете используется бинарное кодирование жанров, что позволяет легко фильтровать фильмы по категориям. Например, фильм может принадлежать к нескольким жанрам одновременно.

\subsection*{Временные метки}
Позволяют анализировать временные закономерности в поведении пользователей, такие как сезонность предпочтений или изменение интересов с течением времени.

\section{Интеграция данных в проект}

\subsection*{Хранение данных в бэкенде}
Данные из датасета MovieLens 100K были интегрированы в базу данных проекта. Для этого использовалась следующая структура:

\begin{itemize}
    \item \textbf{Таблица \texttt{movies}}:
    \begin{itemize}
        \item \texttt{id} (уникальный идентификатор фильма).
        \item \texttt{title} (название фильма).
        \item \texttt{release\_date} (дата выпуска).
        \item \texttt{genres} (список жанров в формате JSON или отдельная таблица для связи многие-ко-многим).
        \item \texttt{imdb\_url} (ссылка на IMDb).
    \end{itemize}
    \item \textbf{Таблица \texttt{users}}:
    \begin{itemize}
        \item \texttt{id} (уникальный идентификатор пользователя).
        \item \texttt{age} (возраст).
        \item \texttt{gender} (пол).
        \item \texttt{occupation} (профессия).
        \item \texttt{zip\_code} (почтовый индекс).
    \end{itemize}
    \item \textbf{Таблица \texttt{ratings}}:
    \begin{itemize}
        \item \texttt{user\_id} (ID пользователя).
        \item \texttt{movie\_id} (ID фильма).
        \item \texttt{rating} (оценка от 1 до 5).
        \item \texttt{timestamp} (временная метка).
    \end{itemize}
    \item \textbf{Таблица \texttt{tags}} (дополнительно):
    \begin{itemize}
        \item \texttt{id} (уникальный идентификатор тега).
        \item \texttt{movie\_id} (ID фильма).
        \item \texttt{tag} (текстовый тег, например, <<классика>>, <<комедия>>, <<для семьи>>).
    \end{itemize}
\end{itemize}

\subsection*{CRUD операции}
Бэкенд поддерживает стандартные CRUD операции:
\begin{itemize}
    \item \textbf{Create}: Добавление новых фильмов, пользователей, оценок и тегов.
    \item \textbf{Read}: Получение информации о фильмах, пользователях, оценках и тегах.
    \item \textbf{Update}: Обновление существующих записей.
    \item \textbf{Delete}: Удаление записей.
\end{itemize}

Пример API-запроса для получения информации о фильме:
\begin{verbatim}
GET /movies/{id}
\end{verbatim}

Пример ответа:
\begin{verbatim}
{
  "id": 1,
  "title": "Toy Story (1995)",
  "release_date": "1995-01-01",
  "genres": ["Animation", "Comedy", "Family"],
  "imdb_url": "http://us.imdb.com/M/title-exact?Toy%20Story%20(1995)"
}
\end{verbatim}

\section{Рекомендательная система (ML)}

\subsection*{Генерация рекомендаций}
Рекомендательная система использует данные из таблиц \texttt{ratings} и \texttt{movies} для генерации персонализированных рекомендаций. Основные подходы:
\begin{itemize}
    \item \textbf{Коллаборативная фильтрация}: Рекомендации на основе схожести пользователей или фильмов. Например, если два пользователя оценили несколько фильмов одинаково, система может рекомендовать фильмы, которые понравились одному из них, другому.
    \item \textbf{Контентная фильтрация}: Рекомендации на основе жанров и тегов фильмов. Например, если пользователь часто смотрит комедии, система будет рекомендовать другие комедии.
    \item \textbf{Гибридные методы}: Комбинация коллаборативной и контентной фильтрации для повышения точности рекомендаций.
\end{itemize}

\subsection*{Обновление рекомендаций}
Рекомендации генерируются раз в сутки с использованием ML-моделей. Результаты сохраняются в отдельную таблицу:
\begin{itemize}
    \item \textbf{Таблица \texttt{recommendations}}:
    \begin{itemize}
        \item \texttt{user\_id} (ID пользователя).
        \item \texttt{movie\_id} (ID рекомендованного фильма).
        \item \texttt{score} (оценка релевантности).
    \end{itemize}
\end{itemize}

Пример API-запроса для получения рекомендаций:
\begin{verbatim}
GET /recommendations/{user_id}
\end{verbatim}

Пример ответа:
\begin{verbatim}
[
  {
    "movie_id": 2,
    "title": "Jumanji (1995)",
    "score": 0.95
  },
  {
    "movie_id": 3,
    "title": "Grumpier Old Men (1995)",
    "score": 0.92
  }
]
\end{verbatim}

\section{Аналитика и оптимизация}

\subsection*{Анализ данных}
\begin{itemize}
    \item \textbf{Распределение оценок}: Анализ того, как пользователи оценивают фильмы. Например, можно выявить, какие фильмы получают самые высокие или самые низкие оценки.
    \item \textbf{Популярность жанров}: Определение самых популярных жанров среди пользователей. Например, комедии могут быть более популярны среди молодых пользователей, а драмы — среди старшей аудитории.
    \item \textbf{Активность пользователей}: Выявление наиболее активных пользователей, которые оставляют больше всего оценок и просматривают больше всего фильмов.
\end{itemize}

\subsection*{Оптимизация}
\begin{itemize}
    \item \textbf{Кэширование}: Для ускорения доступа к часто запрашиваемым данным (например, топ-10 фильмов). Кэширование позволяет снизить нагрузку на базу данных и ускорить время отклика.
    \item \textbf{Индексация}: Оптимизация запросов к базе данных с использованием индексов. Например, индексы могут быть созданы для полей \texttt{user\_id} и \texttt{movie\_id} в таблице \texttt{ratings}.
    \item \textbf{Масштабирование}: Поддержка большого количества пользователей и фильмов за счёт горизонтального масштабирования. Например, можно использовать шардирование базы данных для распределения нагрузки.
\end{itemize}

\section{Дополнительные возможности и улучшения}

\subsection*{Персонализация пользовательского опыта}
\begin{itemize}
    \item \textbf{Рекомендации на основе контекста}: Учёт текущего времени, дня недели или сезона для предоставления более релевантных рекомендаций. Например, в выходные дни можно рекомендовать более длинные фильмы, а в будни — короткометражки.
    \item \textbf{Анализ поведения пользователей}: Использование данных о просмотрах и оценках для создания более точных профилей пользователей. Например, можно учитывать, какие жанры пользователь предпочитает в разное время суток.
\end{itemize}

\subsection*{Интеграция с внешними сервисами}
\begin{itemize}
    \item \textbf{IMDb API}: Использование API IMDb для получения дополнительной информации о фильмах, такой как рейтинги, отзывы и трейлеры.
    \item \textbf{Социальные сети}: Интеграция с социальными сетями для предоставления рекомендаций на основе интересов друзей и знакомых.
\end{itemize}

\subsection*{Улучшение рекомендательной системы}
\begin{itemize}
    \item \textbf{Использование глубокого обучения}: Внедрение нейронных сетей для анализа сложных паттернов в данных и повышения точности рекомендаций.
    \item \textbf{Обратная связь от пользователей}: Сбор и анализ отзывов пользователей для постоянного улучшения рекомендательной системы.
\end{itemize}

Использование датасета MovieLens 100K позволило создать мощную базу для онлайн-кинотеатра, включая хранение данных, CRUD операции и рекомендательную систему. Интеграция ML-моделей обеспечивает персонализированный подход к каждому пользователю, что повышает удовлетворённость и вовлечённость. Проект продолжает развиваться, добавляя новые функции и улучшая существующие. В будущем планируется расширение функционала, включая более глубокий анализ данных, интеграцию с внешними сервисами и использование передовых методов машинного обучения для ещё более точных рекомендаций.

\endinput